\section{Method: roles, dispatch envelope, and gates}\label{sec:method}

\subsection{Role specialization over heterogeneous backbones}
Roles decouple \emph{what a sub-mission needs} from \emph{which model
runs it}. Four roles are defined initially---\textit{architect}
(decomposition, coupling and gate design), \textit{implementer}
(module development and self-test), \textit{verifier} (integration
runs, gate evidence, failure diagnosis), \textit{documenter}
(documentation and paper material)---each bound to a backbone model in
an owner-editable \texttt{roles.json}. The mapping is data-driven: a
seven-model bake-off on an identical goal (\S\ref{sec:ablation})
measures acceptance, tool economy, and wall clock per model, and the
owner ratifies the initial assignment.

\subsection{Agent-driven dispatch under a hard envelope}
A running orchestrator mission dispatches sub-missions through two
bridge tools. \texttt{dispatch\_sub\_mission(goal, repo, role|model,
budget)} queues a child mission; \texttt{await\_children(timeout)}
polls until every child reaches a terminal status (success, failure,
stop, or \emph{awaiting acceptance}---an orchestrator does not block on
the owner's terminal decision) and returns per-child summaries. While
awaiting, the orchestrator's heartbeat thread keeps its own liveness
current, and await time is bounded (default 25\,min, hard cap 60\,min)
so a stuck child cannot silently consume the orchestrator's budget.

Crucially, the agent's platform account is \emph{read-only}: it cannot
create missions through the ordinary endpoint. Instead the dispatch
endpoint authenticates the \emph{parent's claim token} and enforces a
server-side envelope:

\begin{itemize}
\item \textbf{Repository whitelist.} Children may only target the six
project repositories (plus whitelisted co-repositories).
\item \textbf{Model requirement.} Every child must pin a model; a
model-null child could queue forever behind a default daemon busy
running its own orchestrator (deadlock by construction).
\item \textbf{Budget clamps.} Child budgets are clamped to caps
(80 tool calls, 7200\,s wall, $3{\times}10^6$ tokens at time of
writing) regardless of requested values.
\item \textbf{Fan-out and rate limits.} At most 8 children per
orchestrator; at most 24 agent-originated missions per rolling 24\,h.
\item \textbf{Tree depth.} Missions form a forest of depth at most 2:
children cannot dispatch, bounding recursion structurally.
\end{itemize}

The envelope is enforced in one transaction on the dispatch path;
nothing relies on agent cooperation. Children inherit the parent's
organization scope, so the owner's mission list displays the whole tree,
and each dispatch appends events to both parent and child sessions.

\subsection{Cross-repository worktrees}
A mission's worktree set covers its primary repository plus any
co-repositories, checked out sibling under one mission directory; the
agent's workspace root then spans the coupled solvers, while commit and
test tools remain anchored to the primary repository with explicit
\texttt{git -C} for co-repositories. Push is disabled in every
worktree; integration to shared remotes happens only through human
review, preserving a single write path.

\subsection{Machine-checkable stage gates}
Promotion on the coupling ladder is decided exclusively by gates
scripts evaluate: impulse and broken-bond energy ledgers closed to
tolerance (S1), force/moment/work conjugacy between ice--ship reaction
history and structural response (S2), total-force conservation with
temporal alignment for one-way fluid loading (S3), and clock
synchronization with sub-cycling, conservation, rollback, and
checkpoint identities in the closed loop (S4). An agent's summary can
never substitute for a gate; the verifier role assembles gate evidence,
the platform registers runs and artifacts with checksums, and the owner
approves promotion through the same approval chain used for acceptance.
Each stage pre-registers its hypothesis and criteria in an append-only
file before execution.
